Uma \emph{máquina de Turing} é um tipo de mecanismo idealizado de computação.
Ela consiste em uma \emph{fita} infinita em um só sentido, dividida em
quadrados, cada um dos quais pode conter um \emph{símbolo} de um
alfabeto predeterminado. A máquina opera movendo uma
\emph{cabeça de leitura e escrita} ao longo da fita. Ela também pode estar em
um de um número predeterminado de \emph{estados}. As ações da cabeça de
leitura e escrita são determinadas por um conjunto de instruções; cada
instrução está condicionada a que a máquina esteja em um certo estado e leia
um certo símbolo, e especifica qual símbolo a máquina escreverá no
quadrado atual, se ela moverá a cabeça de leitura e escrita um
quadrado para a esquerda, para a direita, ou se permanecerá parada, e
para qual estado ela mudará. Se a fita contém uma certa \emph{entrada},
representada como uma sequência de símbolos na fita, e a máquina é
colocada no estado inicial designado com a cabeça de leitura e escrita lendo
o quadrado mais à esquerda da entrada, o conjunto de instruções determinará,
passo a passo, uma sequência de \emph{configurações} da máquina: conteúdo da
fita, posição da cabeça de leitura e escrita e estado da máquina. Caso a
máquina encontre uma configuração na qual o conjunto de instruções não
contém uma instrução para a combinação atual de símbolo lido/estado,
a máquina \emph{para}, e o conteúdo da fita é a \emph{saída}.

Números podem ser representados muito facilmente como sequências de traços na
fita de uma máquina de Turing. Dizemos que uma função $\Nat \to \Nat$ é
\emph{Turing-computável} se há uma máquina de Turing que, sempre que
iniciada com a representação unária de~$n$ como entrada, eventualmente
para com sua fita contendo a representação unária de~$f(n)$ como
saída. Muitas funções aritméticas conhecidas podem ser mostradas, com
facilidade (ou nem tanto), como Turing-computáveis. Muitos outros modelos de
computação além das máquinas de Turing foram propostos; e sempre se
verificou que as funções aritméticas computáveis neles também são
Turing-computáveis. Isso é visto como um suporte para a
\emph{Tese de Church-Turing}: a de que toda função aritmética que pode
efetivamente ser computada é Turing-computável.
